\documentclass[11pt]{article}
\usepackage[margin=25mm]{geometry}
\usepackage{fontspec}
\setmainfont{Latin Modern Roman}
\usepackage{amsmath,amssymb,mathtools}
\usepackage{unicode-math}
\setmathfont{Latin Modern Math}
\usepackage{xurl,hyperref}
\hypersetup{colorlinks=true,urlcolor=blue}
\setcounter{secnumdepth}{0}
\setlength{\parindent}{0pt}
\setlength{\parskip}{5pt}
\setlength{\emergencystretch}{3em}
\newcommand{\sref}[2]{\hyperlink{source:#1:#2}{[S#2]}}
\newcommand{\cataloguescope}[1]{\paragraph{Recorded target.}#1}
\newcommand{\R}{\mathbb R}
\newcommand{\C}{\mathbb C}
\newcommand{\N}{\mathbb N}
\newcommand{\Z}{\mathbb Z}
\newcommand{\Q}{\mathbb Q}
\newcommand{\F}{\mathbb F}
\newcommand{\eps}{\varepsilon}
\begin{document}
% UnsolvedMath ID 20002539; source catalogue code AIM-PHYSICS-0047
\setcounter{section}{20002538}
\section{Quantitative simplex and dimension bounds for minimum-potential equal-mass configurations}\label{20002539}
{\small UnsolvedMath ID 20002539 · Mathematical Physics}
\subsection{Definitions and mathematical statement}

\paragraph{Shared notation.}$\mathbb N=\{1,2,\ldots\}$; the standard number fields have their usual meanings. All additional conventions are stated locally.


\paragraph{Definitions and notation.}In the Newtonian $n$-body problem, positions $q_i\in\mathbb R^d$ and masses $m_i>0$ obey $m_i\ddot q_i=\nabla_{q_i}U$, where $U(q)=\sum_{i<j}m_im_j/|q_i-q_j|$ and the gravitational constant is normalized to one. Here $T=\tfrac12\sum_i m_i|\dot q_i|^2$, $E=T-U$, and a collision has $q_i=q_j$. The center of mass can be fixed at zero; moment of inertia is $I=\sum_i m_i|q_i|^2$. A central configuration has $\nabla_{q_i}U=-\lambda m_iq_i$ in center-of-mass coordinates for a common scalar $\lambda$. After fixing $I=1$, it is a critical point of $U$ on the collision-free inertia ellipsoid. Congruent configurations and scales are factored out for a finiteness count. A relative equilibrium is a rigidly rotating solution with time-independent shape. For equal masses, a geometric symmetry can combine an orthogonal transformation with a permutation of particle labels, preserving the normalized configuration. Minimizing the positive Newtonian function $U$ on $I=1$ is the source’s convention; physical potential energy is $-U$.

\paragraph{Catalogue wording retained for traceability.}
Problem 10. Is the central configuration for equal masses with minimal potential on the ellipsoid of moment of inertia 1 necessarily symmetric? (Rick Moeckel)
\paragraph{Scope and source qualification.}The source remarks that the symmetry expectation appears false for n=47 and discusses normalization. This historical indication is disclosed; no current universal symmetry theorem or reviewed counterexample is claimed. The dimension and meaning of symmetric must be stated in a proposed theorem. The catalogue title/status was paired with a machine-generated research attempt; that attempt is not peer-reviewed evidence of settlement. The statement covers the original source questions rather than only the title’s subcase.
\subsection{Short English statement}
Determine whether equal-mass minimum-potential central configurations must have nontrivial symmetry.
\subsection{Research attempt}
\paragraph{Model, date and investigated scope.}
OpenAI GPT-6 Astra (configured model \texttt{gpt-6-astra}), ultra reasoning; generated 11 October 2026. Proved the sharp simplex lower bound and symmetric minimizers when d>=n-1. Lower-dimensional equal-mass minimizing symmetry and the source n=47 indication remain unverified.

\paragraph{Sources and literature check.}
Read the supplied complete definition, catalogue wording, scope qualifications and source list for this canonical ID. The original-source citations are inherited as supplied references; no fresh primary-source or novelty search was performed for this attempt. No current open-status claim is inferred from this calculation.

\paragraph{The maximal-dimensional equal-mass case.}Normalize all masses to one, center of mass zero and \(I=\sum_i|q_i|^2=1\). Let \(M=n(n-1)/2\). Expanding squares gives \(\sum_{i<j}|q_i-q_j|^2=n\). The function \(x\mapsto x^{-1/2}\) is strictly convex on positive numbers, so Jensen's inequality yields
\[U=\sum_{i<j}|q_i-q_j|^{-1}\ge M\left(\frac nM\right)^{-1/2}=\frac{M^{3/2}}{\sqrt n}.
\]
Equality holds exactly when every mutual distance is \(\sqrt{2/(n-1)}\). Such configurations are regular simplices and exist when \(d\ge n-1\): take the vectors \(e_i-n^{-1}\mathbf1\) in the sum-zero subspace of \(\mathbb R^n\), rescaled to inertia one. They have all permutation symmetries realized by orthogonal maps.
\paragraph{Why the low-dimensional problem remains.}If \(d<n-1\), n distinct equidistant points cannot exist: their centered Gram matrix has rank \(n-1\), as follows by expanding their equal squared distances. The Jensen bound is then unattainable. It provides no information on whether the actual minimizing configuration has a nontrivial symmetry. A minimizer does exist on the normalized compact sphere because \(U\) diverges at collisions, but compactness does not constrain its stabilizer.
\paragraph{Scope.}This known simplex calculation answers only a dimension range in which the equality geometry is available. It does not verify or reject the source's reported n=47 indication, whose dimension and symmetry convention matter. Proving a symmetry theorem or certifying an asymmetric minimizer in the intended lower-dimensional setting remains the first gap. The physical sign \(-U\) is not confused with the source's minimization of positive \(U\).

\paragraph{Verification and adversarial audit.}
Checked mass/inertia normalization, pair-distance identity, strict convexity, simplex construction, Gram rank and potential sign. This was an LLM adversarial audit, not independent human verification; no theorem prover was used. Prompts A then B applied to this ID: Checked mass/inertia normalization, pair-distance identity, strict convexity, simplex construction, Gram rank and potential sign. Scope remains limited: Determine the symmetry of actual minimizers in dimensions below n-1, where the simplex bound is unattainable.

Coordinator checked the pair-distance identity, Jensen equality geometry and rank obstruction below the simplex dimension. This additional check was model-assisted, not independent human verification.

\paragraph{Tools, process and contributions.}
Python 3 and local shell used to read supplied files and serialize this manuscript; mathematics derived symbolically in the text, with no numerical experiment or formal proof assistant. Contribution prompt A was followed by prompt B. GPT-6 Astra with ultra reasoning drafted the research and performed its stated self-audit; any additional coordinator contribution is identified in the verification and audit. A GPT-6-based Codex coordinator assembled the manuscript and checked repository integration; its exact variant and reasoning configuration were not exposed. No human mathematical contribution, independent human proof check or formal proof certification is supplied. Complete derivations and the consulted source references are included here; private reasoning transcripts are not published.

\paragraph{Final status and first remaining gap.}
\textbf{UNRESOLVED}. Determine the symmetry of actual minimizers in dimensions below n-1, where the simplex bound is unattainable.

\section{Completion Estimate}
\noindent\textbf{COMPLETION: 0\%}
The calculation is an elementary benchmark or conditional reduction; it establishes no claimed new portion of the complete source target.

\subsection{Sources}
\begin{itemize}
\item[\hypertarget{source:20002539:1}{S1}] ULAM.AI and the UnsolvedMath Contributors, UnsolvedMath, record 20002539 (AIM-PHYSICS-0047). Catalogue curation under CC BY 4.0; definitions and English exposition added locally.
\url{https://huggingface.co/datasets/ulamai/UnsolvedMath}.
\item[\hypertarget{source:20002539:2}{S2}] American Institute of Mathematics, Variational Methods in Celestial Mechanics; notes by Kuo-Chang Chen, version 22 June 2003, source item 10.
\url{https://aimath.org/WWN/varcelest/varcelest.pdf}.
\end{itemize}
\label{end:20002539}
\end{document}
